\begin{enumerate}
	\item \textbf{Affirmation 1 : VRAIE.}

	En effet, pour tout $n$ entier naturel non nul, on a :

	$\dfrac{3n^2+4n+7}{6n^2+1}= \dfrac{n^2 \left(3+\dfrac{4}{n}+\dfrac{7}{n^2}\right)}{n^2 \left(6+\dfrac{1}{n^2}\right)}= \dfrac{3+\dfrac{4}{n}+\dfrac{7}{n^2}}{6+\dfrac{1}{n^2}}$

	$\lim\limits_{n\to +\infty} \dfrac{1}{n}=\lim\limits_{n\to +\infty} \dfrac{1}{n^2}=0$

	donc par limite de la somme :\quad $\lim\limits_{n \to +\infty} 3+\dfrac{4}{n}+\dfrac{7}{n^2}=3$\quad et\quad $\lim\limits_{n \to +\infty} 6+\dfrac{1}{n^2}=6$

	et donc par limite du quotient \quad$\lim\limits_{n\to +\infty } \dfrac{3+\dfrac{4}{n}+\dfrac{7}{n^2}}{6+\dfrac{1}{n^2}} = \dfrac{3}{6}=\dfrac{1}{2}$

	D'autre part :\quad $\lim\limits_{n\to +\infty} \dfrac{1}{2}=\dfrac{1}{2}$.

	La suite $\left(u_n\right)$ est donc encadrée par deux suites ayant pour limite $\dfrac{1}{2}$ donc, d'après le théorème des gendarmes, $\lim\limits_{n\to +\infty} u_n=\dfrac{1}{2}$.

	\emph{Remarque :} il ne faut pas se laisser déstabiliser par le fait que tous les termes de la suite sont \textbf{strictement} supérieurs à $\dfrac{1}{2}$, quand on \og passe à la limite\fg{}, les inégalités deviennent large. Par exemple, la suite de terme général $\dfrac{1}{n}$ a des termes tous strictement plus grands que 0, et pourtant, sa limite est 0.


	\item  \textbf{Affirmation 2 : FAUSSE.}

Sur $[-1~;~3]$, la fonction dérivée $h'$ n'est pas croissante, la fonction $h$ n'est donc pas convexe sur $[-1~;~3]$.

	\item \textbf{Affirmation 3 : VRAIE.}

Il y a $10 \times 10 \times 10 \times 10 \times 3 \times 2 = \np{60000}$ codes possibles.  (10 valeurs possibles pour chacun des chiffres et 3 valeurs possibles pour la première lettre et deux valeurs possible pour la deuxième car elles doivent être distinctes.)

Déterminons le nombre de code ne contenant pas de 0 :

Il y a $9\times 9 \times 9 \times 9 \times 3 \times 2 = \np{39366}$ codes ne contenant pas de 0. (9 valeurs possibles pour chacun des chiffres et 3 valeurs possibles pour la première lettre et deux valeurs possible pour la deuxième car elles doivent être distinctes.)

Il y a donc $\np{60000} -\np{39366}=\np{20634}$ codes contenant au moins un zéro.

	\item \textbf{Affirmation 4 : VRAIE.}


$f$ est dérivable sur $]0~;~+\infty[$ comme produit de fonctions dérivables sur cet intervalle et pour tout réel $x$ strictement positif en dérivant ce produit  on a : 

$f'(x) = 1 \times \ln(x) + x\times \dfrac{1}{x}= \ln(x) + 1$, donc

$xf'(x) - f(x)= x \left(\ln(x) + 1\right) - x\ln(x) = x\ln(x) + x - x\ln(x) = x$.

Conclusion : $f$ est bien solution sur $]0~;~+\infty[$ de l'équation différentielle $xy' - y = x$.
\end{enumerate}

\bigskip
